% Generated from locale/tr/content/set-theory/choice/wellorderingproblem.tex; do not hand-edit.


\olfileid[tr]{sth}{choice}{woproblem}
\olsection{İyi Sıralama Sorunu}

İyi Sıralama'yı kabul edip etmememize oldukça çok şeyin bağlı olduğu
açıktır. Fakat bu ilke hakkındaki tartışma, ona denk bir ilkeye, Seçim
Aksiyomu'na odaklanma eğilimindedir. Bu nedenle şimdi dikkatimizi ona
çevireceğiz (ve denkliği ispatlayacağız).

\citeyear{Cantor1883} yılında Cantor, İyi Sıralama Aksiyomu'nu
desteklediğini belirterek onu ``bana temel görünen, sonuçları bakımından
zengin ve genel geçerliliği bakımından özellikle dikkate değer bir düşünce
yasası'' diye adlandırdı (aktaran \citeauthor{Potter2004}
\citeyear[p.~243]{Potter2004}). Fakat Cantor sonunda bu ``Aksiyomun''
ispata gereksinimi olduğuna ikna oldu. Matematik topluluğu da öyle.

Sorun \citeyear{Zermelo1904} yılında Zermelo tarafından ``çözüldü''.
Çözümünü açıklamak için bazı tanımlara ihtiyacımız var.

\begin{defn}
Bir $f$ fonksiyonu, her $x\in\dom{f}$ için $f(x)\in x$ ise ve ancak o
zaman bir \emph{seçim fonksiyonudur}. $f$ bir seçim fonksiyonu ve
$\dom{f}=A\setminus\{\emptyset\}$ ise ve ancak o zaman $f$'ye
\emph{$A$ için bir seçim fonksiyonu} deriz.
\end{defn}

Sezgisel olarak $A$ için bir seçim fonksiyonu, her (boş olmayan)
$x\in A$ kümesinden belirli bir $f(x)$ elemanını \emph{seçer}. Öyleyse
Seçim Aksiyomu şudur:

\begin{axiom}[Seçim]
	Her kümenin bir seçim fonksiyonu vardır.
\end{axiom}

Zermelo, Seçim'in iyi sıralamayı gerektirdiğini ve tersini gösterdi:

\begin{thm}[$\ZF$ içinde]\ollabel{thmwochoice}
İyi Sıralama ile Seçim denktir.
\end{thm}

\begin{proof}
\emph{Soldan sağa.} $A$ bir kümeler kümesi olsun. Birleşim Aksiyomu
gereği $\bigcup A$ vardır; dolayısıyla İyi Sıralama gereği
$\bigcup A$'yı iyi sıralayan bir $<$ bağıntısı vardır. Şimdi
$f(x)=\text{$x$'in $<$-en küçük elemanı}$ olsun. Bu, $A$ için bir seçim
fonksiyonudur.

\emph{Sağdan sola.} $A$ sabit olsun. $A=\emptyset$ ise sonuç açıktır;
bundan sonra $A\neq\emptyset$ varsayın. Seçim gereği
$\Pow{A}\setminus\{\emptyset\}$ için bir $f$ seçim fonksiyonu vardır.
Sonlu Ötesi Özyineleme'yi kullanarak bir fonksiyon tanımlayın:
\begin{align*}
	g(0) &= f(A)\\
	g(\alpha) &=
		\begin{cases}
			\text{dur!{}} &\text{eğer }A = \funimage{g}{\alpha}\\
			f(A \setminus \funimage{g}{\alpha}) & \text{aksi hâlde}\\	
		\end{cases}
\end{align*}
``Dur!'' yönergesi, aksi hâlde daha uzun olacak bir tanımın yalnızca
kısaltmasıdır. Yani $A=\funimage{g}{\alpha}$ ilk kez olduğunda bütün
$\alpha\leq\delta$ için $g(\delta)=A$ olsun. İlk aşamada bunun yalnızca
bir \emph{terim} tanımladığından emin olabiliriz (bkz.
\olref[sth][spine][recursion]{transrecursionschema} sonrasındaki
açıklamalar); fakat gerçekten gereksinim duyduğumuz fonksiyonu elde
ettiğimizi göstereceğiz.

$f$ bir seçim fonksiyonu olduğundan her $\alpha$ için (durmadan önce)
$g(\alpha)=f(A\setminus\funimage{g}{\alpha})\in
A\setminus\funimage{g}{\alpha}$; yani
$g(\alpha)\notin\funimage{g}{\alpha}$ olur. Dolayısıyla
$g(\alpha)=g(\beta)$ ise
$g(\beta)\notin\funimage{g}{\alpha}$, yani $\beta\notin\alpha$ olur;
benzer biçimde $\alpha\notin\beta$. Üçlü Karşılaştırma gereği
$\alpha=\beta$. Öyleyse $g$, durmadan önce !!{injective}.

Sonra gerçekten durduğumuzu, yani $A=g[\alpha]$ olan (en küçük) bir
$\alpha$ sıra sayısının bulunduğunu gözlemleyin. Tersini varsayalım;
o zaman $g$ !!{injective} olduğundan her $\alpha$ sıra sayısı için
$\cardless{\alpha}{A}$ olurdu; bu ise \olref[hartogs]{HartogsLemma} ile
çelişir. Dolayısıyla ayrıca $\ran{g\restriction\alpha}=A$ olur.

Bu olgular bir araya getirildiğinde $g\restriction\alpha$, bir sıra
sayısından $A$'ya giden bir !!a{bijection} olur. Artık
$g\restriction\alpha$, $A$'yı iyi sıralamak için kullanılabilir.
\end{proof}

Dolayısıyla İyi Sıralama ve Seçim birlikte ayakta kalır ya da birlikte
çöker. Fakat soru hâlâ duruyor: ayakta mı kalırlar, çökerler mi?

